\documentclass[12pt, ukrainian]{article}
\usepackage[a4paper]{geometry}
\setlength\parindent{0mm}
\geometry{top=1.0cm,bottom=1.0cm,left=1.3cm,right=1.3cm}
\pagestyle{empty}
\usepackage[T2A]{fontenc}
\usepackage[utf8]{inputenc}
\usepackage[ukrainian]{babel}
\usepackage{graphicx}
\usepackage{caption}
\usepackage{caption}
\DeclareCaptionFormat{nocaption}{}
\captionsetup{format=nocaption,aboveskip=0pt,belowskip=0pt}
\usepackage[Export]{adjustbox}
\adjustboxset{max size={0.9\linewidth}{0.9\paperheight}}
\usepackage{verbatim, listings, clrscode3e, fancyvrb}
\def\code#1{\Verb!#1!}
\begin{document}
%begin
\begin {large}

\textbf{\_\_.11.2025 Контрольна робота, варіант  \No { 10 }}\par


{

\begin{center}
\textbf{Завдання 1}
\end{center}
Нехай int var=13; Один потік збільшив змінну var на 3, а інший збільшив її в три рази.

гонка даних (data race), стан гонки (race condition)

Виберіть усі правильні твердження (якщо такі є).\par
\vspace{2mm}
(\,1\,) Тут може бути тільки гонка даних.%bad

(\,2\,) Тут гарантовано є гонка даних.%bad

(\,3\,) Тут гарантовано нема стану гонки, але може бути гонка даних.%bad

(\,4\,) Тут гарантовано є стан гонки.%good



\begin{center}
\textbf{Завдання 2}
\end{center}
Виберіть усі правильні твердження (якщо такі є).\par
\vspace{2mm}

(\,1\,) Об’єкт future можна скопіювати та роздати копії різним потокам, щоб вони змогли побачити результат.%bad

(\,2\,) Чи запустить функція  async виконання в новому потоці можна задавати політикою.%good

(\,3\,) Значення, що є результатом виконання функції, запущеної через async, зберігається безпосередньо в об’єкті проміза (promise).%bad

(\,4\,) Значення, що є результатом виконання функції, запущеної через async, зберігається безпосередньо в об’єкті ф’ючерса (future).%bad



\begin{center}
\textbf{Завдання 3}
\end{center}
Виберіть усі правильні твердження (якщо такі є).\par
\vspace{2mm}

(\,1\,) Операція над вільною від блокувань структурою даних може виконуватись скільки завгодно довго.%good

(\,2\,) На операціях структури даних без блокувань взаємоблокувань потоків виникнути не може.%bad

(\,3\,) Мова С++20 має стандартну бібліотеку конкурентних структур даних.%bad

(\,4\,) Структура даних без блокувань гарантує завершення виконання довільної операції над цією структурою даних за обмежений час.%bad



\begin{center}
\textbf{Завдання 4}
\end{center}
Виберіть усі правильні твердження (якщо такі є).\par
\vspace{2mm}

(\,1\,) Динамічні змінні не можуть розділятися між потоками.%bad

(\,2\,) Кожен потік програми може мати доступ до динамічних даних програми.%good

(\,3\,) Потік мови С++ не може отримати доступ до даних іншого потоку.%bad

(\,4\,) Потік може мати доступ до автоматичних змінних іншого потоку і це безпечно.%bad



\begin{center}
\textbf{Завдання 5}
\end{center}
У цьому фрагменті коду
\begin{verbatim}
template <class T> void f(T a1, T a2){}

void g(){
    int i=5;
    f(unique_ptr<int>(new int(++i)), unique_ptr<int>(new int(++i)));
\end{verbatim}

(\,1\,) є невизначена поведінка%good

(\,2\,) є стан гонки%bad

(\,3\,) можливий виток ресурсів%bad

(\,4\,) є гонка даних%bad


}
\end {large}


\end{document}

